\documentclass{article}

% ============================================================
% MA 213 Lab Activity Template
% Student-facing worksheet template for lab development.
% ============================================================

\usepackage[margin=1in]{geometry}
\usepackage{xcolor}
\usepackage{parskip}
\usepackage{array}
\usepackage{enumitem}
\usepackage{listings}

\definecolor{codegray}{gray}{0.97}
\definecolor{huddleblue}{RGB}{214,234,255}
\definecolor{predictionyellow}{RGB}{255,242,200}
\definecolor{debatepink}{RGB}{255,225,225}

\lstset{
  basicstyle=\ttfamily\small,
  columns=fullflexible,
  frame=single,
  framerule=0.4pt,
  backgroundcolor=\color{codegray},
  breaklines=true,
  showstringspaces=false,
  keepspaces=true,
  xleftmargin=0pt,
  xrightmargin=0pt
}

\newcommand{\coursenumber}{MA 213}
\newcommand{\weekplaceholder}{1}
\newcommand{\labplaceholder}{1}
\newcommand{\labtitle}{R Intro and First Data Exploration}

\newcommand{\nameblank}{\rule{1.75in}{0.4pt}}
\newcommand{\idblank}{\rule{1.55in}{0.4pt}}
\newcommand{\shortblank}{\rule{1.6in}{0.4pt}}
\newcommand{\longblank}{\rule{0.93\linewidth}{0.4pt}}

\newcommand{\responsebox}[2]{
  \noindent\fbox{%
    \begin{minipage}{0.95\linewidth}
      \textbf{#1} #2
      \vspace{0.55in}
    \end{minipage}
  }
}

\newcommand{\linedbox}[1]{
  \noindent\fbox{%
    \begin{minipage}{0.95\linewidth}
      #1
      \vspace{0.18in}

      \longblank

      \vspace{0.18in}
      \longblank

      \vspace{0.18in}
      \longblank
    \end{minipage}
  }
}

\newcommand{\callout}[3]{
  \noindent\colorbox{#1}{%
    \makebox[\dimexpr0.95\linewidth-2\fboxsep\relax][l]{\textbf{#2}}%
  }

  \noindent\fbox{%
    \begin{minipage}{0.93\linewidth}
      #3
    \end{minipage}
  }
}

\begin{document}

\begin{center}
  {\LARGE \coursenumber{} Week \weekplaceholder{}: Lab \labplaceholder{}}\\
  {\large \labtitle{}}
\end{center}

\begin{enumerate}[label=\arabic*.,leftmargin=*,itemsep=.7em]
  \item \textbf{Your name:} \nameblank \hfill \textbf{BU ID:} \idblank
  \item \textbf{Partner's name:} \nameblank \hfill \textbf{BU ID:} \idblank
\end{enumerate}

\textbf{Big idea:} In this lab, you will practice using R to import and explore data sets, identify variables, and write clear first interpretations.
\textbf{Do not use GenAI tools for this lab.} The point is to practice your own analysis work, coding skills, and statistical reasoning.

\textbf{Submit this worksheet at the end of the lab.}

\textbf{Files used:} \texttt{sleep.csv} and \texttt{burger.csv}

\textbf{Skills practiced:} \texttt{read.csv()}, \texttt{head()}, \texttt{dim()}, \texttt{names()}, \texttt{str()}, \texttt{table()}, variables, and observational units.

\textbf{Your mission:} Work with your partner to get familiar with R/RStudio, import two data sets, and explain what the structure of each data set reveals.

\section*{Icebreaker}

Before opening R: introduce yourself to your partner and answer the warm-up question below.

\responsebox{Warm-up response:}{What is one thing you want your partner to know about you?}

\section*{Getting Started}

Open RStudio. Make sure \texttt{Lab1.R}, \texttt{sleep.csv}, and \texttt{burger.csv} are in the same folder. Then open \texttt{Lab1.R} and use it as your starter script.

\begin{lstlisting}
# Check where R is currently looking for files.
getwd()

# Read in the sleep data.
sleep_data <- read.csv("sleep.csv")

# First look.
head(sleep_data)
\end{lstlisting}

\callout{huddleblue}{HUDDLE UP}{Before calculating anything, what do you think one row represents? What variables look most important for understanding the sleep data?}

\newpage

\section*{Question 1. Import and First Look}

Import \texttt{sleep.csv} into R. Then look at the first few rows and discuss what you notice.

\begin{lstlisting}
# Read the CSV file.
sleep_data <- read.csv("sleep.csv")

# Look at the first few rows.
head(sleep_data)
\end{lstlisting}

\linedbox{\textbf{Write 2-3 sentences describing what you see in the data. What do you think this data set is about?}}

\section*{Question 2. Rows, Variables, and Patients}

Use \texttt{dim()}, \texttt{names()}, and \texttt{unique()} to answer the questions below.

\callout{predictionyellow}{PREDICTION TIME}{Before running the code, predict how many patients are in the data set. Is that the same as the number of rows? Why or why not?}

\begin{lstlisting}
# Dimensions: rows and columns.
dim(sleep_data)

# Variable names.
names(sleep_data)

# Number of unique patients.
length(unique(sleep_data$ID))
\end{lstlisting}

\begin{center}
\renewcommand{\arraystretch}{2.3}
\begin{tabular}{|p{0.34\linewidth}|p{0.54\linewidth}|}
\hline
\textbf{Question} & \textbf{Your answer} \\
\hline
How many rows are in the data set? & \\
\hline
How many variables are in the data set? & \\
\hline
What are the variable names? & \\
\hline
How many patients are in the data set? & \\
\hline
\end{tabular}
\end{center}

\callout{huddleblue}{HUDDLE UP}{If R gives an error such as ``cannot open file,'' it usually means R is looking in a different folder. Use \texttt{getwd()} to see your working directory and then set your working directory to the folder containing the lab files.}

\section*{Question 3. What Does One Row Represent?}

Use \texttt{str()} to inspect the variables. Then decide which variables are numerical and which are categorical.

\begin{lstlisting}
str(sleep_data)
\end{lstlisting}

\begin{center}
\renewcommand{\arraystretch}{2.3}
\begin{tabular}{|p{0.24\linewidth}|p{0.27\linewidth}|p{0.36\linewidth}|}
\hline
\textbf{Variable} & \textbf{Type or role} & \textbf{What it means} \\
\hline
\texttt{extra} & & \\
\hline
\texttt{group} & & \\
\hline
\texttt{ID} & & \\
\hline
\end{tabular}
\end{center}

\linedbox{\textbf{What does one row represent in this data set? Use the variables to justify your answer.}}

\newpage

\section*{Question 4. Burger Choices and Contingency Tables}

Import \texttt{burger.csv}. Then use \texttt{table()} to make a contingency table for burger place and gender.

\begin{lstlisting}
burger_data <- read.csv("burger.csv")

table(burger_data$best_burger_place, burger_data$gender)
\end{lstlisting}

\callout{huddleblue}{HUDDLE UP}{What does each cell in the table count? Does the table answer ``most visited burger place'' directly, or do you need to combine counts across gender?}

\begin{center}
\renewcommand{\arraystretch}{2.25}
\begin{tabular}{|p{0.42\linewidth}|p{0.50\linewidth}|}
\hline
\textbf{Question} & \textbf{Your answer} \\
\hline
What is the most visited burger place overall? & \\
\hline
What comparison did you make? & \\
\hline
What limitation or caution should readers know? & \\
\hline
\end{tabular}
\end{center}

\section*{Question 5. Close the Lab}

Write a short conclusion that uses both data sets. Focus on what the commands helped you learn, not just the numbers you found.

\callout{debatepink}{DEBATE CORNER}{If you and your partner disagree about what a row represents or which burger place is most visited, write both arguments first. Then decide what claim the evidence actually supports.}

\callout{huddleblue}{CLAIM, EVIDENCE, CONTEXT}{%
The main thing I learned is \shortblank.\\[.6em]
My evidence is \shortblank.\\[.6em]
This matters because \shortblank.
}

\linedbox{\textbf{Final response:} state one claim about the data, cite one command or output as evidence, and explain the context.}

\section*{Optional Challenge}

If your team finishes early, create a one-way table that counts only burger places, without splitting by gender.

\begin{lstlisting}
# Optional challenge starter code.
table(burger_data$best_burger_place)
\end{lstlisting}

\end{document}
